\documentclass[12pt, twoside]{article}
\input{../preamble}
\setlength{\parskip}{0.3\baselineskip}

\fancyhead[LE]{\thepage}
\fancyhead[RO]{Markscheme}
\fancyhead[LO]{La Scuola d'Italia / Dr.\ Huson / IB Math 12 \\* 14 September 2026}

\begin{document}
\subsubsection*{1.2 Classwork: Sampling and Histograms --- Markscheme}

\begin{enumerate}[itemsep=0.3cm]

%--- Problem 1: Sampling terminology and survey design [8 marks] ---
\item[1.] [Maximum mark: 8]

A school of 420 students; the principal asks the first 30 students
through the cafeteria door \emph{``Do you agree that the new lunch menu
is a big improvement?''}

\textbf{(a)} Population: all 420 students at the school. \hfill \textbf{A1}

Sample size: 30. \hfill \textbf{A1}

\emph{Accept ``the 420 students'', ``the whole school'', ``all students at
the school''. Do not accept ``the school'' alone, or ``the students in the
cafeteria'' --- the population is the group the principal wants a
conclusion about, not the group she reached.}

\textbf{(b)} Convenience sampling. \hfill \textbf{A1}

\emph{Accept ``convenience sample'' or ``opportunity sampling''. Do not
accept ``random'', ``systematic'' or ``quota''.}

\textbf{(c)} Any one clear reason, for example: only students who eat in
the cafeteria can be asked, so students who bring lunch from home are
never represented. \hfill \textbf{R1}

\emph{Accept any reason of the right kind --- one day only (Tuesday);
only cafeteria users; the first to arrive may be one year group or one
class leaving a lesson early; the sample is not chosen by chance so some
students have no possibility of being asked. R1 requires a stated reason
connected to \emph{who is left out or over-represented}. ``It is too
small'' alone earns nothing --- 30 out of 420 is a workable sample size;
the defect is how it was chosen, not its size.}

\textbf{(d)} The wording invites the student to agree --- it suggests the
answer (``a big improvement'') before the student has given an opinion,
so it pushes replies towards ``yes''. \hfill \textbf{R1}

A neutral rewrite, for example: \emph{``What is your opinion of the new
lunch menu?''} \hfill \textbf{A1}

\emph{Accept any rewrite that offers the options evenhandedly, e.g.
``Do you think the new lunch menu is better, worse, or about the same as
the old one?'' The R1 and the A1 are independent: a candidate who
rewrites the question neutrally without explaining the fault earns the
A1 only.}

\textbf{(e)} Number the students $1$ to $420$ (a list or the school
roll). \hfill \textbf{A1}

Use the GDC's random integer generator to produce 30 different numbers
from 1 to 420, and survey those students. \hfill \textbf{A1}

\emph{\textbf{GDC:} \texttt{menu} $\rightarrow$ \texttt{Probability}
$\rightarrow$ \texttt{Random} $\rightarrow$ \texttt{Integer},
\texttt{randInt(1,420)}.}

\emph{Accept drawing 30 names from a hat, or any method in which every
student has an equal chance of selection and the selection is made by
chance. The first A1 is for giving every student a label; the second is
for a chance mechanism \emph{and} for discarding repeats (``30 different
numbers''). A candidate who only writes ``pick 30 students at random''
earns nothing --- that restates the task without giving a method.}

\newpage

%--- Problem 2: Grouping continuous data and drawing a histogram [7 marks] ---
\item[2.] [Maximum mark: 7]

Heights of 24 tomato plants (raw data on the student sheet), grouped:

{\renewcommand{\arraystretch}{1.3}
\begin{center}
\begin{tabular}{|l|c|c|c|c|c|c|}
\hline
\textbf{Height $h$ (cm)} & $20 \le h < 25$ & $25 \le h < 30$ &
$30 \le h < 35$ & $35 \le h < 40$ & $40 \le h < 45$ & \textbf{Total} \\ \hline
\textbf{Frequency} & 3 & 8 & 5 & 4 & 4 & \textbf{24} \\ \hline
\end{tabular}
\end{center}}

\textbf{(a)} Tally marks recorded for all five classes
\hfill \textbf{A1}

Frequencies $3,\; 8,\; 5,\; 4,\; 4$ and total $24$ all correct
\hfill \textbf{A1}

\emph{The tally A1 is for a tally that is used and that agrees with the
candidate's own frequency column. Accept a completed frequency column
with the tally column left empty for the second A1 only. The three
boundary values $25.0$, $30.0$ and $35.0$
each belong to the class \emph{above} the boundary; a candidate who puts
one of them in the class below produces $3,7,6,4,4$ or similar and does
not earn the second A1.}

\textbf{(b)} Modal class is $25 \le h < 30$. \hfill \textbf{A1}

\emph{Accept ``25 to 30'' or ``the second class''. Do not accept ``8''
alone --- that is the frequency of the modal class, not the class.
FT: award A1 for the class with the largest frequency in the candidate's
own table from (a).}

\textbf{(c)} It is counted in $30 \le h < 35$. \hfill \textbf{A1}

Because the inequality $30 \le h$ includes 30 itself, while the class
below ends at $h < 30$, which excludes it. \hfill \textbf{R1}

\emph{The A1 is for the class, the R1 for the reason.
Accept any wording that identifies the lower boundary as included and the
upper boundary as excluded --- ``the classes are $\le$ on the left and
$<$ on the right'', ``30 is in the class that starts at 30''. The A1 with
no reason earns 1 of 2. A reason resting only on ``it is closer to the
middle'' or ``you round it'' earns nothing.}

\textbf{(d)} Histogram with bars of heights $3,\, 8,\, 5,\, 4,\, 4$ over
the five classes \hfill \textbf{A1}

Bars of equal width drawn between the correct boundaries
$20,25,30,35,40,45$ with no gaps between them \hfill \textbf{A1}

\emph{The first A1 is for the heights matching the candidate's own
frequency column (FT from (a)); the second is for the shape of the
display --- equal widths, bars touching, each bar spanning its own class.
Bars drawn with gaps between them (a bar chart) lose the second A1 only.
Bars centred on the class midpoints rather than
spanning the boundaries also lose the second A1.}

\vspace{-0.35cm}

\begin{center}
\begin{tikzpicture}[x=0.34cm, y=0.34cm]
  \draw[gray!30] (20,0) grid[xstep=1, ystep=1] (45,9);
  \draw[thick, -{Stealth[length=2mm]}] (20,0) -- (46.2,0);
  \node[anchor=west] at (46.6,0) {\footnotesize $h$ (cm)};
  \draw[thick, -{Stealth[length=2mm]}] (20,0) -- (20,9.6);
  \foreach \x in {20,25,30,35,40,45} \draw (\x,0) -- (\x,-0.35) node[below] {\footnotesize \x};
  \foreach \y in {0,2,4,6,8} \draw (20,\y) -- (19.5,\y) node[left] {\footnotesize \y};
  \node[rotate=90, anchor=south] at (16.9,4.5) {\footnotesize frequency};
  \draw[very thick] (20,3) rectangle (25,0);
  \draw[very thick] (25,8) rectangle (30,0);
  \draw[very thick] (30,5) rectangle (35,0);
  \draw[very thick] (35,4) rectangle (40,0);
  \draw[very thick] (40,4) rectangle (45,0);
\end{tikzpicture}
\end{center}

\end{enumerate}
\end{document}
